% HD-F1-T11. GPT-6 Astra Pro, /root/transport; 2026-11-22 UTC.
\section[HD-F1-T11: the forced order-four reflection]{HD-F1-T11: an order-four reflection is forced in the terminal coefficient group}
\label{sec:f1-terminal-order-four}

Every still-missing terminal tensor source of HD-F1-T09 has
\(K=\mathbf Q(i)\), and its actual rank-four coefficient group
\(\Gamma\) contains a reflection of order four.
In particular the terminal alternative over
\(\mathbf Q(\sqrt{-2})\) is absent. The following exact
ramification calculation also gives a second proof of the
exclusion over \(\mathbf Q(\sqrt{-7})\).

\paragraph{1. Original inertia counts and the actual pullback cover.}
Write \(G=\mathfrak S_s\times\Gamma\), \(s=r+1\), and
\(V=\operatorname{Aug}(K^s)\otimes_K M\) as in T09.
The original generating sphere tuple is
\((\pi_j,\gamma_j)\), \(1\le j\le b\), with \(b\ge12\).
Put
\[
 e_j=\operatorname{rank}(\gamma_j-1),\quad
 k=\#\{j:e_j>0\},\quad
 l=\#\{j:e_j\ge2\},\quad E=\sum_j e_j=8+d_0 .
 \tag{F1T11.1}
\]
T09 proves \(1\le d_0\le4\) and \(k+l\le7\) for these
actual inertias. Since each of the \(k-l\) remaining
nontrivial entries has rank one and every other entry has
rank at most four, this gives
\[
            8+d_0=E\le k+3l\le21-2k,\qquad
                      k\le\frac{13-d_0}{2}.
 \tag{F1T11.2}
\]

Let \(C\) be the original full-\(G\) curve and
\(L_0=\mathfrak S_{s-1}\). Retain the actual curves
\[
 X=C/(L_0\times1),\quad
 Z=C/(\mathfrak S_s\times1),\quad
 Y=C/(L_0\times\Gamma).
 \tag{F1T11.3}
\]
Write \(\bar f:Y\to\mathbf P^1\), \(\deg\bar f=s\).
The subgroup intersection and product show that \(X\) is
the normalization of \(Y\times_{\mathbf P^1}Z\):
the intersection is \(L_0\times1\) and the generated
subgroup is \(G\). In particular \(X\to Y\) is the
actual pullback of the full-\(\Gamma\) cover
\(Z\to\mathbf P^1\), with local inertia
\(\gamma_j^{\,\ell_x}\) at a point of ramification
index \(\ell_x\) for \(\bar f\).

The parabolic multiplicity \(d_0\ge1\) implies
\(g(Z)\ge4\). Indeed the complex representations \(M\)
and \(\bar M\) are inequivalent because their character
field is nonreal. Each has multiplicity \(d_0\) in
\(H^1(Z,\mathbf C)\), so that space has dimension at
least \(8d_0\). If \(o_j=\operatorname{ord}(\gamma_j)\)
at the \(k\) nontrivial coefficient branch values, define
\[
 \chi_0=-2+\sum_{e_j>0}(1-o_j^{-1})
                    =\frac{2g(Z)-2}{|\Gamma|}>0 .
 \tag{F1T11.4}
\]
The equality is the ordinary Riemann--Hurwitz formula for
this specified Galois cover.

\paragraph{2. The exact weighted ramification identity.}
Give every other point of \(\mathbf P^1\) the value \(o=1\).
At \(x\in Y\) put
\[
 o'_x=\frac{o_{\bar f(x)}}{\gcd(o_{\bar f(x)},\ell_x)},\qquad
 \chi_Y=2g(Y)-2+\sum_{x\in Y}(1-(o'_x)^{-1}).
 \tag{F1T11.5}
\]
Only finitely many summands are nonzero. Substitution of
Riemann--Hurwitz for the actual degree-\(s\) map
\(\bar f\), and \(\sum_{x\mapsto z}\ell_x=s\), gives
the identity
\[
 \chi_Y=s\chi_0+
  \sum_{x\in Y}
    \frac{\ell_x-\gcd(\ell_x,o_{\bar f(x)})}
         {o_{\bar f(x)}} .
 \tag{F1T11.6}
\]
For \(o=1\), its summand is \(\ell_x-1\).
Every term is nonnegative. At each of the \(b-k\)
original true branch values with \(\gamma_j=1\),
one has \(\pi_j\ne1\). Its ramification contribution
in \(\bar f\) is consequently
\(\sum_{x\mapsto z_j}(\ell_x-1)\ge1\). Thus
\[
                    \chi_Y\ge s\chi_0+b-k
                                 \ge s\chi_0+12-k.
 \tag{F1T11.7}
\]
These are actual extra ramification points of \(\bar f\);
they are not counted as branch points of the pulled-back
\(\Gamma\)-system.

\paragraph{3. Exclusion of groups with only involutory reflections.}
Assume every reflection in \(\Gamma\) has order two.
For any nontrivial \(\gamma\in\Gamma\), write
\(e(\gamma)=\operatorname{rank}(\gamma-1)\) and
\(o(\gamma)=\operatorname{ord}(\gamma)\). If \(e=1\),
then \(1-o^{-1}=1/2=e/2\). If \(e\ge2\), then
\(1-o^{-1}<1\le e/2\). Applying this directly to the
actual local powers \(\gamma_j^{\,\ell_x}\) gives
\[
             \chi_Y\le2g(Y)-2+\frac12\sum_x e'_x .
 \tag{F1T11.8}
\]
The T09 quotient has parabolic dimension \(6+d_0\);
its actual period-rank bound makes \(g(Y)=0\) or \(1\),
as explicitly proved in T10, paragraph 1. Its Euler formula
gives \(\sum e'_x=14+d_0\) at genus zero and
\(\sum e'_x=6+d_0\) at genus one. Hence, uniformly,
\[
                  \chi_Y\le5+\frac{d_0}{2}.
 \tag{F1T11.9}
\]
For genus one the stronger upper bound is
\(3+d_0/2\), but it is not needed.
Combining (F1T11.2), (F1T11.7) and (F1T11.9) yields
\[
 0<s\chi_0
       \le k-7+\frac{d_0}{2}
       \le-\frac12,
 \tag{F1T11.10}
\]
a contradiction. Thus \(\Gamma\) has a reflection whose
order is not two.

For any actual reflection matrix over \(K\), its unique
nontrivial eigenvalue is its determinant and therefore is
a root of unity in \(K\). T08--T10 leave only
\(\mathbf Q(i)\) and \(\mathbf Q(\sqrt{-2})\); the latter
has only \(\{\pm1\}\), whereas the former has \(\mu_4\).
Consequently the field is Gaussian and the required
reflection has order four.

\paragraph{The retained target and framework maps.}
The whole calculation concerns the actual normalized
fibre product (F1T11.3). In its Jacobian the target is
still the image of
\[
 \Pi=e_u(1-s^{-1}f^*f_*),\qquad f:X\to Z,
 \quad f_*f^*=[s],\quad (f^*)^*\Theta_X=s\Theta_Z .
 \tag{F1T11.11}
\]
The original primitive corner is
\(|L_0|^{-1}\alpha^*\Pi\alpha_*
 =(\operatorname{av}_{L_0}-\operatorname{av}_{\mathfrak S_s})e_u\),
where \(\alpha:C\to X\). This is the same Rosati-self-adjoint
projector and the same rank-six norm multiplier as in T09.
Its graph correspondences induce the exact \(K_0\), GRR,
Poincare-biextension and flat variational maps retained
there. No determinant of the rank-\(6+d_0\) containing
factor has been substituted for the target determinant,
and the degree \(s\) on the base summand has not been removed.
